% !TeX root = Javier_TARAZONA_CV_fr.tex
% !TeX program = pdflatex
%-------------------------
% Resume in Latex
% Author : Javier Tarazona
% Offre : BNP Paribas - Stage - AI Research Intern (IT Group, ref S_ITGIDF27_0088) - bwelcome.hr.bnpparibas (Avature)
%------------------------

\documentclass[a4paper,10pt]{article}
\DeclareUnicodeCharacter{00A0}{~}
\DeclareUnicodeCharacter{2013}{\textendash}
\DeclareUnicodeCharacter{2014}{\textemdash}
\DeclareUnicodeCharacter{201C}{``}
\DeclareUnicodeCharacter{201D}{''}
\DeclareUnicodeCharacter{2022}{\textbullet}
\DeclareUnicodeCharacter{25E6}{\textopenbullet}
\usepackage[T1]{fontenc}      % Codificación de salida (acentos y símbolos)
\usepackage[utf8]{inputenc}   % Lee texto en UTF-8
\usepackage[empty]{fullpage}
\usepackage[export]{adjustbox}
\usepackage[pdftex]{hyperref}
\usepackage[usenames,dvipsnames]{color}
\usepackage{array}
\usepackage{comment}
\usepackage{enumitem}         % Control de listas
\usepackage{fancyhdr}
\usepackage{graphicx}
\usepackage{latexsym}
\usepackage{lmodern}
\usepackage{marvosym}
\usepackage{tabularx}
\usepackage{titlesec}
\usepackage{verbatim}
\usepackage{xcolor}

% --- Viñetas limpias y seguras ---
\setlist[itemize]{label=\textbullet, labelsep=0.5em, leftmargin=*}

% ===== CORRECTIFS ATS =====
\input{glyphtounicode}        % chaque glyphe -> vrai texte Unicode (ligatures decomposees)
\pdfgentounicode=1
\usepackage[none]{hyphenat}   % point 5 : plus aucun mot-cle coupe
\sloppy                       %           evite les debordements de marge
\usepackage{textcomp}         % point 6 : apostrophe ASCII au lieu de U+2019
\catcode`\'=\active
\def'{\textquotesingle}
\hypersetup{                  % point 8 : metadonnees indexables
  pdftitle={Javier Andres TARAZONA JIMENEZ - CV - AI Research Intern - Stage de fin d'études en Intelligence Artificielle},
  pdfauthor={Javier Andres TARAZONA JIMENEZ},
  pdfsubject={AI Research Intern - stage de fin d'études Bac+5 de 6 mois à partir d'avril 2027 - recherche en Intelligence Artificielle, Deep Learning, LLM, NLP},
  pdfkeywords={AI Research Intern, Intelligence Artificielle, recherche, Deep Learning, Machine Learning,
               Large Language Models (LLM), NLP, IA générative, fiabilité des LLM, hallucinations,
               apprentissage auto-supervisé, Vision Transformers, Computer Vision, expérimentations,
               évaluation de modèles, Python, PyTorch, scikit-learn, probabilités, statistique, optimisation,
               Bac+5, Diplôme d'Ingénieur, Grande École, stage de fin d'études},
  pdfborder={0 0 0}           % pas de cadre autour des liens (le soulignement bleu suffit)
}
% ==========================

\pagestyle{empty}

% Adjust margins
\addtolength{\oddsidemargin}{-0.45in}
\addtolength{\evensidemargin}{-0.45in}
\addtolength{\textwidth}{1.15in}
\addtolength{\topmargin}{-.62in}
\addtolength{\textheight}{1.20in}

\urlstyle{same}

\raggedbottom
\raggedright
\setlength{\tabcolsep}{0in}
\renewcommand{\arraystretch}{0.98}

%\definecolor{mydarkgreen}{RGB}{0, 140, 123}
\definecolor{mydarkgreen}{RGB}{0, 0, 0}

% Sections formatting
\titleformat{\section}{
  \vspace{-6pt}\scshape\raggedright\large
}{}{0em}{}[\color{mydarkgreen}\titlerule \vspace{-8pt}]

%-------------------------
% Custom commands

% Indented paragraph helper: \indentbox[ancho]{contenido}
% Uso: \indentbox{Texto...}  o  \indentbox[1.5em]{Texto...}
% [t] + \raggedright + \strut : lignes alignees en haut, sans espaces etires ni chevauchement.
\newcommand{\indentbox}[2][1em]{\hspace*{#1}\parbox[t]{\dimexpr\linewidth-#1\relax}{\raggedright #2\strut}}

% Photo et QR code desactives : dans Workday une photo casse l'analyse du texte voisin.


%-------------------------------------------
%%%%%%  CV STARTS HERE  %%%%%%%%%%%%%%%%%%%%%%%%%%%%


\begin{document}

%----------HEADING-----------------
% Nom seul sur la premiere ligne, puis coordonnees en texte brut dans le corps (pas d'en-tete, pas de minipage).

\noindent{\Large\bfseries\textcolor{mydarkgreen}{Javier Andres TARAZONA JIMENEZ}}\par
\vspace{3pt}
\noindent{\footnotesize javitar06@gmail.com \textbullet{} +33 7 46 36 97 75 \textbullet{} Massy, France}\par
\vspace{1pt}
\noindent{\footnotesize\href{https://www.linkedin.com/in/javier-andres-tarazona-jimenez-84b489222/}{\textcolor{blue}{\underline{linkedin.com/in/javier-andres-tarazona-jimenez-84b489222}}} \textbullet{} \href{https://github.com/JavierTarazona06}{\textcolor{blue}{\underline{github.com/JavierTarazona06}}}}\par

\vspace{0.12cm}

{\centering
  {\large\bfseries AI Research Intern - Stage de fin d'études en Intelligence Artificielle}\\[2pt]
  {\footnotesize Élève ingénieur en dernière année (Grande École, Bac+5) \textbullet{} Disponible à partir d'avril 2027 pour 6 mois}\\[3pt]
  {\footnotesize\textit{Recherche en Deep Learning : Large Language Models (LLM) et NLP, Vision Transformers, apprentissage auto-supervisé}}\par}

\vspace{0.05cm}

%-----------FORMATION-----------------
\section{\textcolor{mydarkgreen}{\textbf{FORMATION}}}
{\raggedright
\textbf{ENSTA Paris (Institut Polytechnique de Paris)} \textbullet{} Palaiseau, France \textbullet{} \textit{\footnotesize 2025/09 - 2027/09}\\
{\footnotesize\textit{Diplôme d'Ingénieur (Bac+5) - Master of Science in Engineering, parcours Intelligence Artificielle}}\\[3pt]
{\footnotesize Cours : Deep Learning, Language Modeling, Probability and Machine Learning, Reinforcement Learning, Programmation GPU}\\[4pt]

\textbf{Universidad Nacional de Colombia (UNAL)} \textbullet{} Bogotá, Colombie \textbullet{} \textit{\footnotesize 2021/01 - 2025/07}\\
{\footnotesize\textit{Cursus d'ingénieur - Ingénierie des Systèmes et Informatique}}\\[3pt]
{\footnotesize 12e université d'Amérique latine (QS 2026). Cours : Probabilités et statistique, Modèles stochastiques, Optimisation, Méthodes numériques}\\[4pt]

\textbf{University of Saskatchewan} \textbullet{} Saskatoon, Canada \textbullet{} \textit{\footnotesize 2024/09 - 2024/12}\\
{\footnotesize\textit{Échange académique - Informatique : Algorithmes, Ingénierie logicielle, Systèmes d'exploitation}}\par}

\vspace{0.16cm}
%-----------EXPERIENCE-----------------
\section{\textcolor{mydarkgreen}{\textbf{EXPÉRIENCE PROFESSIONNELLE}}}
{\raggedright
\textbf{ENSTA Paris - Laboratoire U2IS} \textbullet{} Palaiseau, France \textbullet{} \textit{\footnotesize 2026/05 - 2026/08}\\
{\footnotesize\textit{Stagiaire de recherche en Intelligence Artificielle et Vision 3D}}\\[3pt]
\indentbox{\footnotesize \textbf{- Apprentissage auto-supervisé :} pré-entraînement par prédiction de vues (\textbf{PyTorch}, \textbf{Vision Transformers}) pour la structure 3D.}\\
\indentbox{\footnotesize \textbf{- Curriculum learning :} expérimentations de synthèse de nouvelles vues sur des formes 3D synthétiques de complexité croissante.}\\
\indentbox{\footnotesize \textbf{- Données d'entraînement :} pipeline Python/Blender générant des datasets multi-vues (rendu, caméras, métadonnées).}\\[4pt]

\textbf{APEDYS91} \textbullet{} Paris, France \textbullet{} \textit{\footnotesize 2025/10 - 2026/02}\\
{\footnotesize\textit{Ingénieur Machine Learning / NLP}}\\[3pt]
\indentbox{\footnotesize \textbf{- IA générative et NLP :} correction de texte en français par pipeline hybride LanguageTool et \textbf{Large Language Models (LLM)} locaux.}\\
\indentbox{\footnotesize \textbf{- Fiabilité des LLM :} guardrails anti-hallucination (NER spaCy, contrôle des entités et nombres) ; modèle choisi selon RAM/VRAM.}\\[4pt]

\textbf{Engin A.I} \textbullet{} Bogotá, Colombie (à distance, États-Unis) \textbullet{} \textit{\footnotesize 2025/02 - 2025/06}\\
{\footnotesize\textit{Ingénieur Logiciel et Vision par Ordinateur}}\\[3pt]
\indentbox{\footnotesize \textbf{- Vision par ordinateur} (OpenCV, clustering non supervisé) \textbf{:} \textbf{+50\% de vitesse, +60\% de précision} en détection de voies routières.}\\
\indentbox{\footnotesize \textbf{- Analyse de données :} trajectoires de roue (WheelPath) corrélées aux zones de détérioration (+10\% de satisfaction utilisateur).}\\
\indentbox{\footnotesize \textbf{- Qualité du code Python :} refactorisation modulaire de 2 bibliothèques, \textbf{+70\% de score qualité} (maintenabilité, tests).}\\[4pt]

\textbf{Engin A.I} \textbullet{} Bogotá, Colombie (à distance, États-Unis) \textbullet{} \textit{\footnotesize 2024/01 - 2024/07}\\
{\footnotesize\textit{Ingénieur Logiciel et Vision par Ordinateur}}\par}


\vspace{0.16cm}
%-----------PROJETS-----------------
\section{\textcolor{mydarkgreen}{\textbf{PROJETS}}}
{\raggedright
\textbf{"Tameo" - Vision embarquée pour bateau autonome} \textbullet{} \textit{\footnotesize 2025/09 - 2027/06}\\
{\footnotesize\textit{Équipe de 5, membre puis chef de projet depuis 2026/09 - Monaco Energy Boat Challenge, AI Class - Python, PyTorch, YOLO}}\\
\indentbox{\footnotesize \textbf{Évaluation de modèles :} fine-tuning, data augmentation et comparaison de modèles de détection ; mAP et F1 d'environ 80\% à 90\%.}\\[3pt]

\textbf{Segmentation de peau par apprentissage statistique} \textbullet{} {\footnotesize Équipe de 3, 3 mois, rapport en ligne}\\
{\footnotesize\href{https://github.com/JavierTarazona06/Vision_LocalFeatures_BayesianClassification_KMeans}{\textcolor{blue}{\underline{github.com/JavierTarazona06/Vision\_LocalFeatures\_BayesianClassification\_KMeans}}}}\\
\indentbox{\footnotesize \textbf{Étude comparative :} Gaussian Naive Bayes, QDA et K-Means (scikit-learn) sur RGB/HSV/YCrCb pour la segmentation pixel à pixel.}\par}


\vspace{0.16cm}
%-----------COMPETENCES TECHNIQUES-----------------
\section{\textcolor{mydarkgreen}{\textbf{COMPÉTENCES TECHNIQUES}}}
{\footnotesize
\setlength{\parindent}{0pt}\setlength{\parskip}{2pt}
\hangindent=1.4em\hangafter=1 \textbf{Recherche en IA :} Deep Learning, Machine Learning, LLM, NLP, IA générative, apprentissage auto-supervisé (self-supervised learning), Computer Vision, évaluation de modèles\par
\hangindent=1.4em\hangafter=1 \textbf{Frameworks et modèles :} Python, PyTorch, scikit-learn, spaCy, OpenCV, NumPy, Pandas, Vision Transformers, Llama, Mistral, Gemma\par
\hangindent=1.4em\hangafter=1 \textbf{Mathématiques :} probabilités et statistique, modèles stochastiques, optimisation, méthodes numériques, algèbre linéaire\par
\hangindent=1.4em\hangafter=1 \textbf{Ingénierie :} Git/GitHub, Docker, FastAPI, SQL, Google Cloud (GCP), CI/CD, C++/C, Java, Agile (SCRUM)\par
}

\vspace{0.16cm}
%-----------LANGUES-----------------
\section{\textcolor{mydarkgreen}{\textbf{LANGUES}}}
{\footnotesize
\noindent \textbf{Anglais :} courant (C1 - IELTS 2024) \textbullet{} \textbf{Français :} avancé (B2 - DELF 2024) \textbullet{} \textbf{Espagnol :} langue maternelle \textbullet{} \textbf{Portugais :} notions
}

\vspace{0.16cm}
%--------DISTINCTIONS------------
\section{\textcolor{mydarkgreen}{\textbf{DISTINCTIONS}}}
{\footnotesize
\noindent \textbf{Bourse Eiffel Excellence} \textbullet{} bourse d'excellence du gouvernement français pour le Master à ENSTA Paris\\[1pt]
\noindent \textbf{Mention d'Honneur - ICPC Colombie 2023} \textbullet{} XXXVII Marathon National de Programmation ACIS REDIS\\[1pt]
\noindent \textbf{Excellence académique} \textbullet{} UNAL, exonération des frais de scolarité
}


%-------------------------------------------
\end{document}
